\subsection{Arquitectura del Proyecto}

% Ajustes locales para evitar desbordes visuales
\setlength{\tabcolsep}{4pt}
\setlength{\LTleft}{0pt}
\setlength{\LTright}{0pt}
\lstset{basicstyle=\ttfamily\small, breaklines=true, breakatwhitespace=true, columns=fullflexible, keepspaces=true}

\section{Patrón Arquitectónico: Feature-First}

El proyecto implementa una arquitectura \textbf{Feature-First} (también conocida como arquitectura por módulos funcionales) \cite{feature_first_architecture}.

\subsection{Definición}

En una arquitectura Feature-First, el código se organiza por funcionalidades del negocio en lugar de por capas técnicas (MVC, MVVM) \cite{clean_architecture}. Cada feature contiene todas las capas necesarias (UI, lógica, datos) en un módulo autocontenido.

\subsection{Justificación}

\begin{enumerate}
    \item \textbf{Escalabilidad modular:} Nuevos features se agregan sin modificar código existente
    \item \textbf{Mantenibilidad:} Cambios en un feature no afectan otros (bajo acoplamiento)
    \item \textbf{Claridad:} Fácil localizar funcionalidades específicas
    \item \textbf{Desarrollo paralelo:} 2 desarrolladores pueden trabajar en features diferentes sin conflictos de merge
    \item \textbf{Testing aislado:} Cada feature puede tener sus propios tests unitarios
\end{enumerate}

\section{Estructura de Carpetas Completa}

\begin{lstlisting}[caption={Arbol del proyecto}]
lib/
|-- main.dart
|-- core/
|   |-- constants/
|   |   `-- app_colors.dart
|   `-- models/
|       `-- finance_models.dart
`-- features/                          # 11 modulos funcionales
    |-- auth/
    |   |-- screens/
    |   |   |-- login_screen.dart
    |   |   `-- register_screen.dart
    |   `-- services/
    |       `-- auth_service.dart
    |-- onboarding/                    # Bienvenida
    |   `-- screens/
    |       |-- onboarding_screen.dart
    |       |-- initial_setup_1_screen.dart
    |       |-- initial_setup_2_screen.dart
    |       `-- initial_setup_3_screen.dart
    |-- home/                          # Dashboard
    |   `-- screens/
    |       |-- main_screen.dart
    |       `-- home_screen.dart
    |-- scan/                          # OCR
    |   `-- screens/
    |       |-- scan_ticket_screen.dart
    |       |-- scan_preview_screen.dart
    |       |-- scan_processing_screen.dart
    |       |-- scan_result_screen.dart
    |       `-- view_ticket_screen.dart
    |-- transactions/                  # Transacciones
    |   `-- screens/
    |       |-- transactions_screen.dart
    |       |-- add_transaction_screen.dart
    |       |-- register_expense_screen.dart
    |       `-- register_income_screen.dart
    |-- budgets/                       # Presupuestos
    |   `-- screens/
    |       |-- budgets_screen.dart
    |       |-- new_budget_screen.dart
    |       |-- edit_budget_screen.dart
    |       `-- edit_incomes_screen.dart
    |-- debts/                         # Deudas
    |   `-- screens/
    |       |-- debts_screen.dart
    |       |-- new_debt_screen.dart
    |       `-- edit_debt_screen.dart
    |-- categories/                    # Categorias
    |   `-- screens/
    |       `-- new_category_screen.dart
    |-- notifications/                 # Notificaciones
    |   `-- screens/
    |       `-- notifications_screen.dart
    |-- profile/                       # Perfil
    |   `-- screens/
    |       |-- profile_screen.dart
    |       `-- profile_settings_screen.dart
    `-- explore/                       # Explorar
        `-- screens/
            `-- explore_screen.dart   # Placeholder
\end{lstlisting}

\section{Descripción de Módulos}

\subsection{Core (Compartido)}

\begin{description}
    \item[app\_colors.dart:] Paleta de colores Material Design 3 (Primary, Secondary, Error, etc.)
    \item[finance\_models.dart:] Modelos de datos: \texttt{User}, \texttt{Transaction}, \texttt{Budget}, \texttt{Debt}, \texttt{Income}
\end{description}

\subsection{Features (11 Módulos Funcionales)}

\begin{longtable}{|p{2.8cm}|p{1.8cm}|p{7.6cm}|}
\caption{Features implementados} \\
\hline
\textbf{Feature} & \textbf{Pantallas} & \textbf{Responsabilidad} \\
\hline
\endfirsthead
\hline
\textbf{Feature} & \textbf{Pantallas} & \textbf{Responsabilidad} \\
\hline
\endhead

Auth & 2 & Login, registro, logout, gestión de sesión con JWT \\
\hline
Onboarding & 4 & Tutorial inicial (5 slides) + 3 setups de configuración \\
\hline
Home & 2 & Dashboard principal con gráficas de pastel, Bottom Navigation \\
\hline
Scan (OCR) & 5 & \textbf{Feature estrella:} Captura, procesamiento OCR, extracción de datos \\
\hline
Transactions & 4 & Historial con gráficas de barras, registro manual de ingresos/gastos \\
\hline
Budgets & 4 & CRUD de presupuestos, edición de ingresos, metodología Dave Ramsey \\
\hline
Debts & 3 & CRUD de deudas, planificación de pagos, alertas de vencimiento \\
\hline
Categories & 1 & Creación de categorías personalizadas (máx. 10) \\
\hline
Notifications & 1 & Historial de notificaciones, configuración de alertas \\
\hline
Profile & 2 & Visualización y edición de perfil (nombre, email, password) \\
\hline
Explore & 1 & Placeholder para futuras funcionalidades \\
\hline
\end{longtable}

\section{Navegación y Rutas}

\subsection{Configuración de Named Routes}

El proyecto utiliza \textbf{Named Routes} de Flutter \cite{flutter_navigation} para navegación estructurada.

\begin{lstlisting}[caption=Configuración en main.dart]
MaterialApp(
  title: 'FINARA',
  theme: ThemeData(useMaterial3: true),
  initialRoute: '/login',
  routes: {
    '/login': (context) => const LoginScreen(),
    '/register': (context) => const RegisterScreen(),
    '/onboarding': (context) => const OnboardingScreen(),
    '/initial-setup-1': (context) => const InitialSetup1Screen(),
    '/initial-setup-2': (context) => const InitialSetup2Screen(),
    '/initial-setup-3': (context) => const InitialSetup3Screen(),
    '/main': (context) => const MainScreen(),
    '/register-expense': (context) => const RegisterExpenseScreen(),
    '/register-income': (context) => const RegisterIncomeScreen(),
    '/edit-incomes': (context) => const EditIncomesScreen(),
    '/new-budget': (context) => const NewBudgetScreen(),
    '/new-debt': (context) => const NewDebtScreen(),
    '/new-category': (context) => const NewCategoryScreen(),
    '/notifications': (context) => const NotificationsScreen(),
    '/profile-settings': (context) => const ProfileSettingsScreen(),
    '/scan-ticket': (context) => const ScanTicketScreen(),
    '/scan-preview': (context) => const ScanPreviewScreen(),
    '/scan-processing': (context) => const ScanProcessingScreen(),
    '/scan-result': (context) => const ScanResultScreen(),
    '/edit-debt': (context) => const EditDebtScreen(),
    '/edit-budget': (context) => const EditBudgetScreen(),
    '/view-ticket': (context) => const ViewTicketScreen(),
  },
)
\end{lstlisting}

\textbf{Total de rutas configuradas:} 23

\subsection{Flujo Principal de Navegación}

\begin{table}[H]
\centering
\begin{tabularx}{\textwidth}{|l|X|}
\hline
\textbf{Paso} & \textbf{Descripcion} \\
\hline
1 & Login o Register \\
\hline
2 & Onboarding para usuario nuevo \\
\hline
3 & Setup 1 \\
\hline
4 & Setup 2 \\
\hline
5 & Setup 3 \\
\hline
6 & Main Screen \\
\hline
\end{tabularx}
\caption{Flujo de navegación principal (nuevo usuario)}
\label{fig:flujo_navegacion}
\end{table}

\subsection{Bottom Navigation (Main Screen)}

\begin{table}[H]
\centering
\begin{tabularx}{\textwidth}{|>{\centering\arraybackslash}p{2.4cm}|X|}
\hline
\textbf{Tab} & \textbf{Pantalla} \\
\hline
Home & \texttt{home\_screen.dart} \\
\hline
Transactions & \texttt{transactions\_screen.dart} \\
\hline
Budgets & \texttt{budgets\_screen.dart} \\
\hline
Debts & \texttt{debts\_screen.dart} \\
\hline
Explore & \texttt{explore\_screen.dart} \\
\hline
\end{tabularx}
\caption{Configuración de Bottom Navigation}
\end{table}

\section{Patrones de Diseño Implementados}

\begin{description}
    \item[StatefulWidget:] Usado en todas las pantallas con estado mutable (formularios, listas dinámicas)
    \item[Singleton:] \texttt{AuthService} implementa singleton para gestión única de sesión
    \item[Factory Pattern:] Modelos con constructores factory para deserialización JSON
    \item[Builder Pattern:] ListView.builder para listas eficientes con lazy loading
    \item[Observer:] setState() para notificar cambios de estado a la UI
\end{description}
